% ─────────────────────────────────────────────────────────────────────────────
% preamble.tex — Préambule commun à tous les fichiers maîtres du projet.
%
% Usage dans chaque fichier maître :
%   \documentclass[12pt,twoside,openright]{memoir}
%   % optionnel : \def\AVECNOTES{} et/ou \def\AVECILLUSTRATIONS{}
%   % ─────────────────────────────────────────────────────────────────────────────
% preamble.tex — Préambule commun à tous les fichiers maîtres du projet.
%
% Usage dans chaque fichier maître :
%   \documentclass[12pt,twoside,openright]{memoir}
%   % optionnel : \def\AVECNOTES{} et/ou \def\AVECILLUSTRATIONS{}
%   % ─────────────────────────────────────────────────────────────────────────────
% preamble.tex — Préambule commun à tous les fichiers maîtres du projet.
%
% Usage dans chaque fichier maître :
%   \documentclass[12pt,twoside,openright]{memoir}
%   % optionnel : \def\AVECNOTES{} et/ou \def\AVECILLUSTRATIONS{}
%   % ─────────────────────────────────────────────────────────────────────────────
% preamble.tex — Préambule commun à tous les fichiers maîtres du projet.
%
% Usage dans chaque fichier maître :
%   \documentclass[12pt,twoside,openright]{memoir}
%   % optionnel : \def\AVECNOTES{} et/ou \def\AVECILLUSTRATIONS{}
%   \input{preamble}
%   \begin{document}
%   ...
% ─────────────────────────────────────────────────────────────────────────────

% ─── PDF VERSION ──────────────────────────────────────────────────────────────
\ifdefined\pdfvariable\pdfvariable minorversion 7\fi

% ─── PACKAGES ─────────────────────────────────────────────────────────────────
\usepackage{fontspec}
\setmainfont{EB Garamond}
\newfontfamily{\arabicfont}{Geeza Pro}
\newcommand{\arab}[1]{{\arabicfont\textdir TRT #1}}
\newfontfamily{\cjkfont}{Hiragino Mincho ProN}
\newcommand{\cjk}[1]{{\cjkfont #1}}
\usepackage[russian,french]{babel}
\usepackage[program=/usr/local/bin/lilypond]{lyluatex}
\usepackage{microtype}
\usepackage{fmtcount}
\usepackage{catchfile}
\usepackage{graphicx}
\usepackage{xcolor}
\usepackage{amsmath}
\usepackage{xparse}
\usepackage{tikz}
\usepackage{afterpage}
\usepackage{needspace}
\usepackage{pdfpages}
\usepackage[absolute,overlay]{textpos}
\usepackage{tabularx}
\usepackage{longtable}
\usepackage{booktabs}
\usepackage{calc}
\usepackage{wrapfig}
\usepackage{hyperref}
\hypersetup{pdfpagelayout=TwoPageRight, hidelinks}

% ─── VERSION / GIT ────────────────────────────────────────────────────────────
\input{version.tex}
\IfFileExists{tete_de_veau_ravigote.gitinfo}%
  {\CatchFileDef{\gitcommit}{tete_de_veau_ravigote.gitinfo}{\endlinechar=-1}}%
  {\def\gitcommit{}}

% ─── GRAYSCALE IMAGE FILTER ───────────────────────────────────────────────────
\graphicspath{{iconography/}}
\newcommand{\figcaptionname}{\footnotesize}
\newcommand{\figcaptiondetail}{\scriptsize}
\newcommand{\bwimage}[2][]{%
  \begin{tikzpicture}
    \node[inner sep=0, outer sep=0pt, fill=white] (img) {\includegraphics[#1]{#2}};
    \begin{scope}[blend mode=saturation]
      \fill[black] (img.south west) rectangle (img.north east);
    \end{scope}
  \end{tikzpicture}%
}

% ─── ESPACEMENT FLOTTANTS ─────────────────────────────────────────────────────
\setlength{\textfloatsep}{8pt}   % espace entre flottant top/bottom et le texte (défaut ~20pt)

% ─── ICONOGRAPHIE ─────────────────────────────────────────────────────────────
% Compiler avec : lualatex "\def\AVECILLUSTRATIONS{}\input{...}"
\ifdefined\AVECILLUSTRATIONS
  \newcommand{\iconographieimg}[2]{%
    \afterpage{%
      \clearpage
      \thispagestyle{empty}%
      \vspace*{\fill}%
      \begin{center}%
        \bwimage[height=0.75\textheight,width=\textwidth,keepaspectratio]{#1}\\[0.6em]%
        {\small\itshape #2}%
      \end{center}%
      \vspace*{\fill}%
      \clearpage
    }%
  }
  \newcommand{\iconographietex}[1]{%
    \afterpage{%
      \thispagestyle{empty}%
      \input{#1}%
      \clearpage
    }%
  }
  % Image en wrapfigure : #1=pos  #2=largeur  #3=code-image  #4=légende
  \newcommand{\iconographiewrapfig}[4]{%
    \setlength{\intextsep}{0pt}      % espace vertical au-dessus/dessous                                                                                                                                 
    \setlength{\columnsep}{8pt}%
    \begin{wrapfigure}{#1}{#2}%
      % \vspace{0.5ex}%
      \centering
      #3%
      \ifx&#4&\else\\[0.2em]{\scriptsize\itshape #4}\\[0.5\baselineskip]\fi%
    \end{wrapfigure}%
  }
  % Bloc centré : [#1]=placement float optionnel (t/b), #2=code-image, #3=légende
  \newcommand{\iconographieinlineblock}[3][]{%
    \ifx&#1&%
      % Pas de placement → inline
      \vspace{0.5\onelineskip}%
      {\centering#2\par}%
      \ifx&#3&\else{\centering\scriptsize\itshape#3\par\vspace{0.5\baselineskip}}\fi%
      \vspace{-0.5\onelineskip}%
    \else
      % Placement fourni → figure flottante
      \begin{figure}[#1]
        \centering#2%
        \ifx&#3&\else\par{\scriptsize\itshape#3}\fi%
      \end{figure}%
    \fi
  }
  % Double image côte à côte :
  %   #1 = placement (bottom/top, défaut: bottom)
  %   #2 = hauteur commune (ex: 5cm)
  %   #3 = image gauche   #5 = image droite
  \NewDocumentCommand{\iconographiedouble}{O{bottom} O{5cm} m O{} m O{} O{}}{%
    \ifthenelse{\equal{#1}{bottom}}{\vspace*{\fill}}{\vspace*{0pt}}%
    \vspace{0.5\onelineskip}%
    \noindent\bwimage[height=#2]{#3}\hfill\bwimage[height=#2]{#5}\par%
    \ifx&#7&%
      \vspace{0.5\onelineskip}%
    \else%
      \vspace{0.2\onelineskip}%
      {\centering\scriptsize\itshape#7\par}%
      \vspace{0.3\onelineskip}%
    \fi%
    \ifthenelse{\equal{#1}{top}}{\vspace*{\fill}}{}%
  }
\else
  \newcommand{\iconographieimg}[2]{}
  \newcommand{\iconographietex}[1]{}
  \newcommand{\iconographiewrapfig}[4]{}
  \NewDocumentCommand{\iconographieinlineblock}{O{} m m}{}
  \NewDocumentCommand{\iconographiedouble}{O{bottom} O{5cm} m O{} m O{} O{}}{}
\fi

% ─── NOTES FINALES ────────────────────────────────────────────────────────────
% Compiler avec : lualatex "\def\AVECNOTES{}\input{...}"
\usepackage{endnotes}
\ifdefined\AVECNOTES
  \newcommand{\nf}[1]{\endnote{#1}}
\else
  \newcommand{\nf}[1]{}
\fi
% Supprimer le titre automatique du package (on gère le nôtre)
\renewcommand{\notesname}{}
% Style des notes : même fonte, corps légèrement réduit
\renewcommand{\enotesize}{\small}
% Format des notes en fin de livre : sans retrait, numéro en corps normal
\makeatletter
\renewcommand\enoteformat{%
  \rightskip\z@ \leftskip\z@ \parindent=0pt
  \leavevmode\theenmark.\enspace}
\makeatother
% Source dans les notes : saut de ligne avant le lien Wikipedia
\newcommand{\source}[1]{\newline\textit{Source~:}~{\upshape\def\_{\textunderscore\allowbreak}\def\%{\char"25\relax\allowbreak}#1}}
% Image seule dans une note (centrée sous le texte)
\newcommand{\nfimg}[2]{\newline\includegraphics[width=#1\linewidth]{#2}}
% Image côte à côte avec le texte : #1=texte, #2=frac image, #3=path, #4=url source
\newcommand{\nfimgblock}[4]{%
  \begin{minipage}[t]{0.45\linewidth}\small#1\end{minipage}%
  \hfill%
  \begin{minipage}[t]{#2\linewidth}\includegraphics[width=\linewidth]{#3}\end{minipage}%
  \source{#4}}
% Filet ornemental (séparateur de scène intra-chapitre)
\newcommand{\scenbreak}{%
  \vspace{\onelineskip}%
  {\centering\rule{0.4\textwidth}{0.4pt}\par}%
  \vspace{\onelineskip}%
}
% Titre de chapitre dans le bloc de notes
\newcommand{\chapnoteheader}[1]{{\normalfont\scshape #1}}
% Sous-chapitre intermédiaire : vide les notes courantes, ouvre un nouveau chapitre, repart à 1
% Usage : \startnewchapter        (titre auto = numéro romain du chapitre suivant)
%         \startnewchapter[Titre] (titre explicite)
% Actif uniquement si \CHAPITRESSEPARES est défini dans le document racine.
\ifdefined\CHAPITRESSEPARES
  \newcommand{\startnewchapter}[1][]{%
    \ifdefined\AVECNOTES
      \cleardoublepage
      \def\@snc@title{#1}%
      {\centering\normalfont\scshape\large
        Notes\ifx\@snc@title\empty\else~\textemdash~#1\fi\par}%
      \vspace{\onelineskip}%
      \begingroup\parindent=0pt
      \theendnotes
      \endgroup
    \fi
    \chapter{\ifx&#1&\Roman{chapter}\else#1\fi}%
    \beginchapternotes{\ifx&#1&\Roman{chapter}\else#1\fi}%
  }
\else
  \newcommand{\startnewchapter}[1][]{}%
\fi
% Notes par chapitre : remet le compteur à 0 pour que chaque acte reparte de 1
\newcommand{\beginchapternotes}[1]{%
  \ifdefined\AVECNOTES
    \setcounter{endnote}{0}%
    \ifdefined\NOTESPARTCHAPITRE\else
      \addtoendnotes{\protect\par\protect\medskip\protect\noindent\protect\chapnoteheader{#1}\protect\par}%
    \fi
  \fi
}

% ─── PAGE GEOMETRY (155×235 mm – Medium 1) ───────────────────────────────────
\setstocksize{235mm}{155mm}
\settrimmedsize{235mm}{155mm}{*}
\settrims{0mm}{0mm}
\setlrmarginsandblock{23mm}{18mm}{*} % inner/outer
\setulmarginsandblock{25mm}{28mm}{*} % top/bottom
\setheadfoot{11mm}{13mm}
\setheaderspaces{*}{7mm}{*}
\checkandfixthelayout

% ─── TYPOGRAPHY ───────────────────────────────────────────────────────────────
\raggedbottom
\renewcommand{\baselinestretch}{1.2}
\setlength{\parindent}{4mm}
\setlength{\parskip}{0pt}
\setlength{\afterchapskip}{0.4\onelineskip}

% ─── DIALOGUE ─────────────────────────────────────────────────────────────────
\newenvironment{dialogue}
  {\par\vspace{0.5\baselineskip}\setlength{\parindent}{0pt}}
  {\par\vspace{0.5\baselineskip}}


% ─── TABLE DES MATIÈRES ───────────────────────────────────────────────────────
\setcounter{tocdepth}{0}
\renewcommand{\chapternumberline}[1]{}           % suppress arabic number, keep Roman title only
\renewcommand{\cftchapterdotsep}{\cftdotsep}     % dot leaders
\renewcommand{\cftchapterfont}{\normalfont\scshape}     % chapter entry: small caps, not bold
\renewcommand{\cftchapterpagefont}{\normalfont\scshape} % page number: small caps, not bold
% TOC title — gallimard chapter style suppresses \printchaptertitle, so we override here
\renewcommand{\printtoctitle}[1]{{\centering\normalfont\scshape\large #1\par}}
\renewcommand{\aftertoctitle}{\par\vspace{\onelineskip}}
% \mytableofcontents : cleardoublepage + suppress self-reference entry
\newcommand{\mytableofcontents}{%
  \cleardoublepage
  \begingroup
  \small\renewcommand{\baselinestretch}{1.0}\selectfont
  \addtocontents{toc}{\protect\setcounter{tocdepth}{-1}}%
  \tableofcontents
  \addtocontents{toc}{\protect\setcounter{tocdepth}{0}}%
  \endgroup
}

% ─── CHAPTER STYLE ────────────────────────────────────────────────────────────
\makeatletter
\makechapterstyle{gallimard}{%
  \renewcommand{\chapnamefont}{\normalfont\scshape\centering}
  \renewcommand{\chapnumfont}{\normalfont\scshape\centering}
  \renewcommand{\chaptitlefont}{\normalfont\scshape\centering}
  \renewcommand{\printchaptername}{}
  \renewcommand{\chapternamenum}{}
  \renewcommand{\printchapternum}{%
    \chapnumfont\Roman{chapter}%
    \markboth{\Roman{chapter}}{\Roman{chapter}}}
  \renewcommand{\afterchapternum}{\par\vskip 0.4\onelineskip}
  \renewcommand{\printchaptertitle}[1]{}
}
\makeatother
\chapterstyle{gallimard}

% ─── RUNNING HEADERS : folio seul, sans auteur ni titre ───────────────────────
\makepagestyle{gallimard}
\makeevenhead{gallimard}{}{\footnotesize\scshape\leftmark}{}
\makeoddhead{gallimard}{}{\footnotesize\scshape\rightmark}{}
\makeevenfoot{gallimard}{}{\thepage}{}
\makeoddfoot{gallimard}{}{\thepage}{}
\makepsmarks{gallimard}{%
  \nouppercaseheads
}
\pagestyle{gallimard}
\aliaspagestyle{chapter}{empty}

% ─── EDITORIAL NOTE (Avertissement, Postfaces) ────────────────────────────────
% 10.95pt (~small), interligne ×1.1, sans alinéa, titre en petites capitales.
\newenvironment{editorialnote}[1]{%
  \begingroup
  \renewcommand{\baselinestretch}{1.1}%
  \small
  \setlength{\parindent}{0pt}%
  \setlength{\parskip}{\onelineskip}%
  {\centering\normalfont\scshape\large #1\par}%
  \vspace{\onelineskip}%
  \noindent\ignorespaces
}{%
  \par\endgroup
}

%   \begin{document}
%   ...
% ─────────────────────────────────────────────────────────────────────────────

% ─── PDF VERSION ──────────────────────────────────────────────────────────────
\ifdefined\pdfvariable\pdfvariable minorversion 7\fi

% ─── PACKAGES ─────────────────────────────────────────────────────────────────
\usepackage{fontspec}
\setmainfont{EB Garamond}
\newfontfamily{\arabicfont}{Geeza Pro}
\newcommand{\arab}[1]{{\arabicfont\textdir TRT #1}}
\newfontfamily{\cjkfont}{Hiragino Mincho ProN}
\newcommand{\cjk}[1]{{\cjkfont #1}}
\usepackage[russian,french]{babel}
\usepackage[program=/usr/local/bin/lilypond]{lyluatex}
\usepackage{microtype}
\usepackage{fmtcount}
\usepackage{catchfile}
\usepackage{graphicx}
\usepackage{xcolor}
\usepackage{amsmath}
\usepackage{xparse}
\usepackage{tikz}
\usepackage{afterpage}
\usepackage{needspace}
\usepackage{pdfpages}
\usepackage[absolute,overlay]{textpos}
\usepackage{tabularx}
\usepackage{longtable}
\usepackage{booktabs}
\usepackage{calc}
\usepackage{wrapfig}
\usepackage{hyperref}
\hypersetup{pdfpagelayout=TwoPageRight, hidelinks}

% ─── VERSION / GIT ────────────────────────────────────────────────────────────
\newcommand{\docversion}{0.5.0}
\newcommand{\docversionordinal}{Deuxième}

\IfFileExists{tete_de_veau_ravigote.gitinfo}%
  {\CatchFileDef{\gitcommit}{tete_de_veau_ravigote.gitinfo}{\endlinechar=-1}}%
  {\def\gitcommit{}}

% ─── GRAYSCALE IMAGE FILTER ───────────────────────────────────────────────────
\graphicspath{{iconography/}}
\newcommand{\figcaptionname}{\footnotesize}
\newcommand{\figcaptiondetail}{\scriptsize}
\newcommand{\bwimage}[2][]{%
  \begin{tikzpicture}
    \node[inner sep=0, outer sep=0pt, fill=white] (img) {\includegraphics[#1]{#2}};
    \begin{scope}[blend mode=saturation]
      \fill[black] (img.south west) rectangle (img.north east);
    \end{scope}
  \end{tikzpicture}%
}

% ─── ESPACEMENT FLOTTANTS ─────────────────────────────────────────────────────
\setlength{\textfloatsep}{8pt}   % espace entre flottant top/bottom et le texte (défaut ~20pt)

% ─── ICONOGRAPHIE ─────────────────────────────────────────────────────────────
% Compiler avec : lualatex "\def\AVECILLUSTRATIONS{}\input{...}"
\ifdefined\AVECILLUSTRATIONS
  \newcommand{\iconographieimg}[2]{%
    \afterpage{%
      \clearpage
      \thispagestyle{empty}%
      \vspace*{\fill}%
      \begin{center}%
        \bwimage[height=0.75\textheight,width=\textwidth,keepaspectratio]{#1}\\[0.6em]%
        {\small\itshape #2}%
      \end{center}%
      \vspace*{\fill}%
      \clearpage
    }%
  }
  \newcommand{\iconographietex}[1]{%
    \afterpage{%
      \thispagestyle{empty}%
      \input{#1}%
      \clearpage
    }%
  }
  % Image en wrapfigure : #1=pos  #2=largeur  #3=code-image  #4=légende
  \newcommand{\iconographiewrapfig}[4]{%
    \setlength{\intextsep}{0pt}      % espace vertical au-dessus/dessous                                                                                                                                 
    \setlength{\columnsep}{8pt}%
    \begin{wrapfigure}{#1}{#2}%
      % \vspace{0.5ex}%
      \centering
      #3%
      \ifx&#4&\else\\[0.2em]{\scriptsize\itshape #4}\\[0.5\baselineskip]\fi%
    \end{wrapfigure}%
  }
  % Bloc centré : [#1]=placement float optionnel (t/b), #2=code-image, #3=légende
  \newcommand{\iconographieinlineblock}[3][]{%
    \ifx&#1&%
      % Pas de placement → inline
      \vspace{0.5\onelineskip}%
      {\centering#2\par}%
      \ifx&#3&\else{\centering\scriptsize\itshape#3\par\vspace{0.5\baselineskip}}\fi%
      \vspace{-0.5\onelineskip}%
    \else
      % Placement fourni → figure flottante
      \begin{figure}[#1]
        \centering#2%
        \ifx&#3&\else\par{\scriptsize\itshape#3}\fi%
      \end{figure}%
    \fi
  }
  % Double image côte à côte :
  %   #1 = placement (bottom/top, défaut: bottom)
  %   #2 = hauteur commune (ex: 5cm)
  %   #3 = image gauche   #5 = image droite
  \NewDocumentCommand{\iconographiedouble}{O{bottom} O{5cm} m O{} m O{} O{}}{%
    \ifthenelse{\equal{#1}{bottom}}{\vspace*{\fill}}{\vspace*{0pt}}%
    \vspace{0.5\onelineskip}%
    \noindent\bwimage[height=#2]{#3}\hfill\bwimage[height=#2]{#5}\par%
    \ifx&#7&%
      \vspace{0.5\onelineskip}%
    \else%
      \vspace{0.2\onelineskip}%
      {\centering\scriptsize\itshape#7\par}%
      \vspace{0.3\onelineskip}%
    \fi%
    \ifthenelse{\equal{#1}{top}}{\vspace*{\fill}}{}%
  }
\else
  \newcommand{\iconographieimg}[2]{}
  \newcommand{\iconographietex}[1]{}
  \newcommand{\iconographiewrapfig}[4]{}
  \NewDocumentCommand{\iconographieinlineblock}{O{} m m}{}
  \NewDocumentCommand{\iconographiedouble}{O{bottom} O{5cm} m O{} m O{} O{}}{}
\fi

% ─── NOTES FINALES ────────────────────────────────────────────────────────────
% Compiler avec : lualatex "\def\AVECNOTES{}\input{...}"
\usepackage{endnotes}
\ifdefined\AVECNOTES
  \newcommand{\nf}[1]{\endnote{#1}}
\else
  \newcommand{\nf}[1]{}
\fi
% Supprimer le titre automatique du package (on gère le nôtre)
\renewcommand{\notesname}{}
% Style des notes : même fonte, corps légèrement réduit
\renewcommand{\enotesize}{\small}
% Format des notes en fin de livre : sans retrait, numéro en corps normal
\makeatletter
\renewcommand\enoteformat{%
  \rightskip\z@ \leftskip\z@ \parindent=0pt
  \leavevmode\theenmark.\enspace}
\makeatother
% Source dans les notes : saut de ligne avant le lien Wikipedia
\newcommand{\source}[1]{\newline\textit{Source~:}~{\upshape\def\_{\textunderscore\allowbreak}\def\%{\char"25\relax\allowbreak}#1}}
% Image seule dans une note (centrée sous le texte)
\newcommand{\nfimg}[2]{\newline\includegraphics[width=#1\linewidth]{#2}}
% Image côte à côte avec le texte : #1=texte, #2=frac image, #3=path, #4=url source
\newcommand{\nfimgblock}[4]{%
  \begin{minipage}[t]{0.45\linewidth}\small#1\end{minipage}%
  \hfill%
  \begin{minipage}[t]{#2\linewidth}\includegraphics[width=\linewidth]{#3}\end{minipage}%
  \source{#4}}
% Filet ornemental (séparateur de scène intra-chapitre)
\newcommand{\scenbreak}{%
  \vspace{\onelineskip}%
  {\centering\rule{0.4\textwidth}{0.4pt}\par}%
  \vspace{\onelineskip}%
}
% Titre de chapitre dans le bloc de notes
\newcommand{\chapnoteheader}[1]{{\normalfont\scshape #1}}
% Sous-chapitre intermédiaire : vide les notes courantes, ouvre un nouveau chapitre, repart à 1
% Usage : \startnewchapter        (titre auto = numéro romain du chapitre suivant)
%         \startnewchapter[Titre] (titre explicite)
% Actif uniquement si \CHAPITRESSEPARES est défini dans le document racine.
\ifdefined\CHAPITRESSEPARES
  \newcommand{\startnewchapter}[1][]{%
    \ifdefined\AVECNOTES
      \cleardoublepage
      \def\@snc@title{#1}%
      {\centering\normalfont\scshape\large
        Notes\ifx\@snc@title\empty\else~\textemdash~#1\fi\par}%
      \vspace{\onelineskip}%
      \begingroup\parindent=0pt
      \theendnotes
      \endgroup
    \fi
    \chapter{\ifx&#1&\Roman{chapter}\else#1\fi}%
    \beginchapternotes{\ifx&#1&\Roman{chapter}\else#1\fi}%
  }
\else
  \newcommand{\startnewchapter}[1][]{}%
\fi
% Notes par chapitre : remet le compteur à 0 pour que chaque acte reparte de 1
\newcommand{\beginchapternotes}[1]{%
  \ifdefined\AVECNOTES
    \setcounter{endnote}{0}%
    \ifdefined\NOTESPARTCHAPITRE\else
      \addtoendnotes{\protect\par\protect\medskip\protect\noindent\protect\chapnoteheader{#1}\protect\par}%
    \fi
  \fi
}

% ─── PAGE GEOMETRY (155×235 mm – Medium 1) ───────────────────────────────────
\setstocksize{235mm}{155mm}
\settrimmedsize{235mm}{155mm}{*}
\settrims{0mm}{0mm}
\setlrmarginsandblock{23mm}{18mm}{*} % inner/outer
\setulmarginsandblock{25mm}{28mm}{*} % top/bottom
\setheadfoot{11mm}{13mm}
\setheaderspaces{*}{7mm}{*}
\checkandfixthelayout

% ─── TYPOGRAPHY ───────────────────────────────────────────────────────────────
\raggedbottom
\renewcommand{\baselinestretch}{1.2}
\setlength{\parindent}{4mm}
\setlength{\parskip}{0pt}
\setlength{\afterchapskip}{0.4\onelineskip}

% ─── DIALOGUE ─────────────────────────────────────────────────────────────────
\newenvironment{dialogue}
  {\par\vspace{0.5\baselineskip}\setlength{\parindent}{0pt}}
  {\par\vspace{0.5\baselineskip}}


% ─── TABLE DES MATIÈRES ───────────────────────────────────────────────────────
\setcounter{tocdepth}{0}
\renewcommand{\chapternumberline}[1]{}           % suppress arabic number, keep Roman title only
\renewcommand{\cftchapterdotsep}{\cftdotsep}     % dot leaders
\renewcommand{\cftchapterfont}{\normalfont\scshape}     % chapter entry: small caps, not bold
\renewcommand{\cftchapterpagefont}{\normalfont\scshape} % page number: small caps, not bold
% TOC title — gallimard chapter style suppresses \printchaptertitle, so we override here
\renewcommand{\printtoctitle}[1]{{\centering\normalfont\scshape\large #1\par}}
\renewcommand{\aftertoctitle}{\par\vspace{\onelineskip}}
% \mytableofcontents : cleardoublepage + suppress self-reference entry
\newcommand{\mytableofcontents}{%
  \cleardoublepage
  \begingroup
  \small\renewcommand{\baselinestretch}{1.0}\selectfont
  \addtocontents{toc}{\protect\setcounter{tocdepth}{-1}}%
  \tableofcontents
  \addtocontents{toc}{\protect\setcounter{tocdepth}{0}}%
  \endgroup
}

% ─── CHAPTER STYLE ────────────────────────────────────────────────────────────
\makeatletter
\makechapterstyle{gallimard}{%
  \renewcommand{\chapnamefont}{\normalfont\scshape\centering}
  \renewcommand{\chapnumfont}{\normalfont\scshape\centering}
  \renewcommand{\chaptitlefont}{\normalfont\scshape\centering}
  \renewcommand{\printchaptername}{}
  \renewcommand{\chapternamenum}{}
  \renewcommand{\printchapternum}{%
    \chapnumfont\Roman{chapter}%
    \markboth{\Roman{chapter}}{\Roman{chapter}}}
  \renewcommand{\afterchapternum}{\par\vskip 0.4\onelineskip}
  \renewcommand{\printchaptertitle}[1]{}
}
\makeatother
\chapterstyle{gallimard}

% ─── RUNNING HEADERS : folio seul, sans auteur ni titre ───────────────────────
\makepagestyle{gallimard}
\makeevenhead{gallimard}{}{\footnotesize\scshape\leftmark}{}
\makeoddhead{gallimard}{}{\footnotesize\scshape\rightmark}{}
\makeevenfoot{gallimard}{}{\thepage}{}
\makeoddfoot{gallimard}{}{\thepage}{}
\makepsmarks{gallimard}{%
  \nouppercaseheads
}
\pagestyle{gallimard}
\aliaspagestyle{chapter}{empty}

% ─── EDITORIAL NOTE (Avertissement, Postfaces) ────────────────────────────────
% 10.95pt (~small), interligne ×1.1, sans alinéa, titre en petites capitales.
\newenvironment{editorialnote}[1]{%
  \begingroup
  \renewcommand{\baselinestretch}{1.1}%
  \small
  \setlength{\parindent}{0pt}%
  \setlength{\parskip}{\onelineskip}%
  {\centering\normalfont\scshape\large #1\par}%
  \vspace{\onelineskip}%
  \noindent\ignorespaces
}{%
  \par\endgroup
}

%   \begin{document}
%   ...
% ─────────────────────────────────────────────────────────────────────────────

% ─── PDF VERSION ──────────────────────────────────────────────────────────────
\ifdefined\pdfvariable\pdfvariable minorversion 7\fi

% ─── PACKAGES ─────────────────────────────────────────────────────────────────
\usepackage{fontspec}
\setmainfont{EB Garamond}
\newfontfamily{\arabicfont}{Geeza Pro}
\newcommand{\arab}[1]{{\arabicfont\textdir TRT #1}}
\newfontfamily{\cjkfont}{Hiragino Mincho ProN}
\newcommand{\cjk}[1]{{\cjkfont #1}}
\usepackage[russian,french]{babel}
\usepackage[program=/usr/local/bin/lilypond]{lyluatex}
\usepackage{microtype}
\usepackage{fmtcount}
\usepackage{catchfile}
\usepackage{graphicx}
\usepackage{xcolor}
\usepackage{amsmath}
\usepackage{xparse}
\usepackage{tikz}
\usepackage{afterpage}
\usepackage{needspace}
\usepackage{pdfpages}
\usepackage[absolute,overlay]{textpos}
\usepackage{tabularx}
\usepackage{longtable}
\usepackage{booktabs}
\usepackage{calc}
\usepackage{wrapfig}
\usepackage{hyperref}
\hypersetup{pdfpagelayout=TwoPageRight, hidelinks}

% ─── VERSION / GIT ────────────────────────────────────────────────────────────
\newcommand{\docversion}{0.5.0}
\newcommand{\docversionordinal}{Deuxième}

\IfFileExists{tete_de_veau_ravigote.gitinfo}%
  {\CatchFileDef{\gitcommit}{tete_de_veau_ravigote.gitinfo}{\endlinechar=-1}}%
  {\def\gitcommit{}}

% ─── GRAYSCALE IMAGE FILTER ───────────────────────────────────────────────────
\graphicspath{{iconography/}}
\newcommand{\figcaptionname}{\footnotesize}
\newcommand{\figcaptiondetail}{\scriptsize}
\newcommand{\bwimage}[2][]{%
  \begin{tikzpicture}
    \node[inner sep=0, outer sep=0pt, fill=white] (img) {\includegraphics[#1]{#2}};
    \begin{scope}[blend mode=saturation]
      \fill[black] (img.south west) rectangle (img.north east);
    \end{scope}
  \end{tikzpicture}%
}

% ─── ESPACEMENT FLOTTANTS ─────────────────────────────────────────────────────
\setlength{\textfloatsep}{8pt}   % espace entre flottant top/bottom et le texte (défaut ~20pt)

% ─── ICONOGRAPHIE ─────────────────────────────────────────────────────────────
% Compiler avec : lualatex "\def\AVECILLUSTRATIONS{}\input{...}"
\ifdefined\AVECILLUSTRATIONS
  \newcommand{\iconographieimg}[2]{%
    \afterpage{%
      \clearpage
      \thispagestyle{empty}%
      \vspace*{\fill}%
      \begin{center}%
        \bwimage[height=0.75\textheight,width=\textwidth,keepaspectratio]{#1}\\[0.6em]%
        {\small\itshape #2}%
      \end{center}%
      \vspace*{\fill}%
      \clearpage
    }%
  }
  \newcommand{\iconographietex}[1]{%
    \afterpage{%
      \thispagestyle{empty}%
      \input{#1}%
      \clearpage
    }%
  }
  % Image en wrapfigure : #1=pos  #2=largeur  #3=code-image  #4=légende
  \newcommand{\iconographiewrapfig}[4]{%
    \setlength{\intextsep}{0pt}      % espace vertical au-dessus/dessous                                                                                                                                 
    \setlength{\columnsep}{8pt}%
    \begin{wrapfigure}{#1}{#2}%
      % \vspace{0.5ex}%
      \centering
      #3%
      \ifx&#4&\else\\[0.2em]{\scriptsize\itshape #4}\\[0.5\baselineskip]\fi%
    \end{wrapfigure}%
  }
  % Bloc centré : [#1]=placement float optionnel (t/b), #2=code-image, #3=légende
  \newcommand{\iconographieinlineblock}[3][]{%
    \ifx&#1&%
      % Pas de placement → inline
      \vspace{0.5\onelineskip}%
      {\centering#2\par}%
      \ifx&#3&\else{\centering\scriptsize\itshape#3\par\vspace{0.5\baselineskip}}\fi%
      \vspace{-0.5\onelineskip}%
    \else
      % Placement fourni → figure flottante
      \begin{figure}[#1]
        \centering#2%
        \ifx&#3&\else\par{\scriptsize\itshape#3}\fi%
      \end{figure}%
    \fi
  }
  % Double image côte à côte :
  %   #1 = placement (bottom/top, défaut: bottom)
  %   #2 = hauteur commune (ex: 5cm)
  %   #3 = image gauche   #5 = image droite
  \NewDocumentCommand{\iconographiedouble}{O{bottom} O{5cm} m O{} m O{} O{}}{%
    \ifthenelse{\equal{#1}{bottom}}{\vspace*{\fill}}{\vspace*{0pt}}%
    \vspace{0.5\onelineskip}%
    \noindent\bwimage[height=#2]{#3}\hfill\bwimage[height=#2]{#5}\par%
    \ifx&#7&%
      \vspace{0.5\onelineskip}%
    \else%
      \vspace{0.2\onelineskip}%
      {\centering\scriptsize\itshape#7\par}%
      \vspace{0.3\onelineskip}%
    \fi%
    \ifthenelse{\equal{#1}{top}}{\vspace*{\fill}}{}%
  }
\else
  \newcommand{\iconographieimg}[2]{}
  \newcommand{\iconographietex}[1]{}
  \newcommand{\iconographiewrapfig}[4]{}
  \NewDocumentCommand{\iconographieinlineblock}{O{} m m}{}
  \NewDocumentCommand{\iconographiedouble}{O{bottom} O{5cm} m O{} m O{} O{}}{}
\fi

% ─── NOTES FINALES ────────────────────────────────────────────────────────────
% Compiler avec : lualatex "\def\AVECNOTES{}\input{...}"
\usepackage{endnotes}
\ifdefined\AVECNOTES
  \newcommand{\nf}[1]{\endnote{#1}}
\else
  \newcommand{\nf}[1]{}
\fi
% Supprimer le titre automatique du package (on gère le nôtre)
\renewcommand{\notesname}{}
% Style des notes : même fonte, corps légèrement réduit
\renewcommand{\enotesize}{\small}
% Format des notes en fin de livre : sans retrait, numéro en corps normal
\makeatletter
\renewcommand\enoteformat{%
  \rightskip\z@ \leftskip\z@ \parindent=0pt
  \leavevmode\theenmark.\enspace}
\makeatother
% Source dans les notes : saut de ligne avant le lien Wikipedia
\newcommand{\source}[1]{\newline\textit{Source~:}~{\upshape\def\_{\textunderscore\allowbreak}\def\%{\char"25\relax\allowbreak}#1}}
% Image seule dans une note (centrée sous le texte)
\newcommand{\nfimg}[2]{\newline\includegraphics[width=#1\linewidth]{#2}}
% Image côte à côte avec le texte : #1=texte, #2=frac image, #3=path, #4=url source
\newcommand{\nfimgblock}[4]{%
  \begin{minipage}[t]{0.45\linewidth}\small#1\end{minipage}%
  \hfill%
  \begin{minipage}[t]{#2\linewidth}\includegraphics[width=\linewidth]{#3}\end{minipage}%
  \source{#4}}
% Filet ornemental (séparateur de scène intra-chapitre)
\newcommand{\scenbreak}{%
  \vspace{\onelineskip}%
  {\centering\rule{0.4\textwidth}{0.4pt}\par}%
  \vspace{\onelineskip}%
}
% Titre de chapitre dans le bloc de notes
\newcommand{\chapnoteheader}[1]{{\normalfont\scshape #1}}
% Sous-chapitre intermédiaire : vide les notes courantes, ouvre un nouveau chapitre, repart à 1
% Usage : \startnewchapter        (titre auto = numéro romain du chapitre suivant)
%         \startnewchapter[Titre] (titre explicite)
% Actif uniquement si \CHAPITRESSEPARES est défini dans le document racine.
\ifdefined\CHAPITRESSEPARES
  \newcommand{\startnewchapter}[1][]{%
    \ifdefined\AVECNOTES
      \cleardoublepage
      \def\@snc@title{#1}%
      {\centering\normalfont\scshape\large
        Notes\ifx\@snc@title\empty\else~\textemdash~#1\fi\par}%
      \vspace{\onelineskip}%
      \begingroup\parindent=0pt
      \theendnotes
      \endgroup
    \fi
    \chapter{\ifx&#1&\Roman{chapter}\else#1\fi}%
    \beginchapternotes{\ifx&#1&\Roman{chapter}\else#1\fi}%
  }
\else
  \newcommand{\startnewchapter}[1][]{}%
\fi
% Notes par chapitre : remet le compteur à 0 pour que chaque acte reparte de 1
\newcommand{\beginchapternotes}[1]{%
  \ifdefined\AVECNOTES
    \setcounter{endnote}{0}%
    \ifdefined\NOTESPARTCHAPITRE\else
      \addtoendnotes{\protect\par\protect\medskip\protect\noindent\protect\chapnoteheader{#1}\protect\par}%
    \fi
  \fi
}

% ─── PAGE GEOMETRY (155×235 mm – Medium 1) ───────────────────────────────────
\setstocksize{235mm}{155mm}
\settrimmedsize{235mm}{155mm}{*}
\settrims{0mm}{0mm}
\setlrmarginsandblock{23mm}{18mm}{*} % inner/outer
\setulmarginsandblock{25mm}{28mm}{*} % top/bottom
\setheadfoot{11mm}{13mm}
\setheaderspaces{*}{7mm}{*}
\checkandfixthelayout

% ─── TYPOGRAPHY ───────────────────────────────────────────────────────────────
\raggedbottom
\renewcommand{\baselinestretch}{1.2}
\setlength{\parindent}{4mm}
\setlength{\parskip}{0pt}
\setlength{\afterchapskip}{0.4\onelineskip}

% ─── DIALOGUE ─────────────────────────────────────────────────────────────────
\newenvironment{dialogue}
  {\par\vspace{0.5\baselineskip}\setlength{\parindent}{0pt}}
  {\par\vspace{0.5\baselineskip}}


% ─── TABLE DES MATIÈRES ───────────────────────────────────────────────────────
\setcounter{tocdepth}{0}
\renewcommand{\chapternumberline}[1]{}           % suppress arabic number, keep Roman title only
\renewcommand{\cftchapterdotsep}{\cftdotsep}     % dot leaders
\renewcommand{\cftchapterfont}{\normalfont\scshape}     % chapter entry: small caps, not bold
\renewcommand{\cftchapterpagefont}{\normalfont\scshape} % page number: small caps, not bold
% TOC title — gallimard chapter style suppresses \printchaptertitle, so we override here
\renewcommand{\printtoctitle}[1]{{\centering\normalfont\scshape\large #1\par}}
\renewcommand{\aftertoctitle}{\par\vspace{\onelineskip}}
% \mytableofcontents : cleardoublepage + suppress self-reference entry
\newcommand{\mytableofcontents}{%
  \cleardoublepage
  \begingroup
  \small\renewcommand{\baselinestretch}{1.0}\selectfont
  \addtocontents{toc}{\protect\setcounter{tocdepth}{-1}}%
  \tableofcontents
  \addtocontents{toc}{\protect\setcounter{tocdepth}{0}}%
  \endgroup
}

% ─── CHAPTER STYLE ────────────────────────────────────────────────────────────
\makeatletter
\makechapterstyle{gallimard}{%
  \renewcommand{\chapnamefont}{\normalfont\scshape\centering}
  \renewcommand{\chapnumfont}{\normalfont\scshape\centering}
  \renewcommand{\chaptitlefont}{\normalfont\scshape\centering}
  \renewcommand{\printchaptername}{}
  \renewcommand{\chapternamenum}{}
  \renewcommand{\printchapternum}{%
    \chapnumfont\Roman{chapter}%
    \markboth{\Roman{chapter}}{\Roman{chapter}}}
  \renewcommand{\afterchapternum}{\par\vskip 0.4\onelineskip}
  \renewcommand{\printchaptertitle}[1]{}
}
\makeatother
\chapterstyle{gallimard}

% ─── RUNNING HEADERS : folio seul, sans auteur ni titre ───────────────────────
\makepagestyle{gallimard}
\makeevenhead{gallimard}{}{\footnotesize\scshape\leftmark}{}
\makeoddhead{gallimard}{}{\footnotesize\scshape\rightmark}{}
\makeevenfoot{gallimard}{}{\thepage}{}
\makeoddfoot{gallimard}{}{\thepage}{}
\makepsmarks{gallimard}{%
  \nouppercaseheads
}
\pagestyle{gallimard}
\aliaspagestyle{chapter}{empty}

% ─── EDITORIAL NOTE (Avertissement, Postfaces) ────────────────────────────────
% 10.95pt (~small), interligne ×1.1, sans alinéa, titre en petites capitales.
\newenvironment{editorialnote}[1]{%
  \begingroup
  \renewcommand{\baselinestretch}{1.1}%
  \small
  \setlength{\parindent}{0pt}%
  \setlength{\parskip}{\onelineskip}%
  {\centering\normalfont\scshape\large #1\par}%
  \vspace{\onelineskip}%
  \noindent\ignorespaces
}{%
  \par\endgroup
}

%   \begin{document}
%   ...
% ─────────────────────────────────────────────────────────────────────────────

% ─── PDF VERSION ──────────────────────────────────────────────────────────────
\ifdefined\pdfvariable\pdfvariable minorversion 7\fi

% ─── PACKAGES ─────────────────────────────────────────────────────────────────
\usepackage{fontspec}
\setmainfont{EB Garamond}
\newfontfamily{\arabicfont}{Geeza Pro}
\newcommand{\arab}[1]{{\arabicfont\textdir TRT #1}}
\newfontfamily{\cjkfont}{Hiragino Mincho ProN}
\newcommand{\cjk}[1]{{\cjkfont #1}}
\usepackage[russian,french]{babel}
\usepackage[program=/usr/local/bin/lilypond]{lyluatex}
\usepackage{microtype}
\usepackage{fmtcount}
\usepackage{catchfile}
\usepackage{graphicx}
\usepackage{xcolor}
\usepackage{amsmath}
\usepackage{xparse}
\usepackage{tikz}
\usepackage{afterpage}
\usepackage{needspace}
\usepackage{pdfpages}
\usepackage[absolute,overlay]{textpos}
\usepackage{tabularx}
\usepackage{longtable}
\usepackage{booktabs}
\usepackage{calc}
\usepackage{wrapfig}
\usepackage{hyperref}
\hypersetup{pdfpagelayout=TwoPageRight, hidelinks}

% ─── VERSION / GIT ────────────────────────────────────────────────────────────
\newcommand{\docversion}{0.5.0}
\newcommand{\docversionordinal}{Deuxième}

\IfFileExists{tete_de_veau_ravigote.gitinfo}%
  {\CatchFileDef{\gitcommit}{tete_de_veau_ravigote.gitinfo}{\endlinechar=-1}}%
  {\def\gitcommit{}}

% ─── GRAYSCALE IMAGE FILTER ───────────────────────────────────────────────────
\graphicspath{{iconography/}}
\newcommand{\figcaptionname}{\footnotesize}
\newcommand{\figcaptiondetail}{\scriptsize}
\newcommand{\bwimage}[2][]{%
  \begin{tikzpicture}
    \node[inner sep=0, outer sep=0pt, fill=white] (img) {\includegraphics[#1]{#2}};
    \begin{scope}[blend mode=saturation]
      \fill[black] (img.south west) rectangle (img.north east);
    \end{scope}
  \end{tikzpicture}%
}

% ─── ESPACEMENT FLOTTANTS ─────────────────────────────────────────────────────
\setlength{\textfloatsep}{8pt}   % espace entre flottant top/bottom et le texte (défaut ~20pt)

% ─── ICONOGRAPHIE ─────────────────────────────────────────────────────────────
% Compiler avec : lualatex "\def\AVECILLUSTRATIONS{}\input{...}"
\ifdefined\AVECILLUSTRATIONS
  \newcommand{\iconographieimg}[2]{%
    \afterpage{%
      \clearpage
      \thispagestyle{empty}%
      \vspace*{\fill}%
      \begin{center}%
        \bwimage[height=0.75\textheight,width=\textwidth,keepaspectratio]{#1}\\[0.6em]%
        {\small\itshape #2}%
      \end{center}%
      \vspace*{\fill}%
      \clearpage
    }%
  }
  \newcommand{\iconographietex}[1]{%
    \afterpage{%
      \thispagestyle{empty}%
      \input{#1}%
      \clearpage
    }%
  }
  % Image en wrapfigure : #1=pos  #2=largeur  #3=code-image  #4=légende
  \newcommand{\iconographiewrapfig}[4]{%
    \setlength{\intextsep}{0pt}      % espace vertical au-dessus/dessous                                                                                                                                 
    \setlength{\columnsep}{8pt}%
    \begin{wrapfigure}{#1}{#2}%
      % \vspace{0.5ex}%
      \centering
      #3%
      \ifx&#4&\else\\[0.2em]{\scriptsize\itshape #4}\\[0.5\baselineskip]\fi%
    \end{wrapfigure}%
  }
  % Bloc centré : [#1]=placement float optionnel (t/b), #2=code-image, #3=légende
  \newcommand{\iconographieinlineblock}[3][]{%
    \ifx&#1&%
      % Pas de placement → inline
      \vspace{0.5\onelineskip}%
      {\centering#2\par}%
      \ifx&#3&\else{\centering\scriptsize\itshape#3\par\vspace{0.5\baselineskip}}\fi%
      \vspace{-0.5\onelineskip}%
    \else
      % Placement fourni → figure flottante
      \begin{figure}[#1]
        \centering#2%
        \ifx&#3&\else\par{\scriptsize\itshape#3}\fi%
      \end{figure}%
    \fi
  }
  % Double image côte à côte :
  %   #1 = placement (bottom/top, défaut: bottom)
  %   #2 = hauteur commune (ex: 5cm)
  %   #3 = image gauche   #5 = image droite
  \NewDocumentCommand{\iconographiedouble}{O{bottom} O{5cm} m O{} m O{} O{}}{%
    \ifthenelse{\equal{#1}{bottom}}{\vspace*{\fill}}{\vspace*{0pt}}%
    \vspace{0.5\onelineskip}%
    \noindent\bwimage[height=#2]{#3}\hfill\bwimage[height=#2]{#5}\par%
    \ifx&#7&%
      \vspace{0.5\onelineskip}%
    \else%
      \vspace{0.2\onelineskip}%
      {\centering\scriptsize\itshape#7\par}%
      \vspace{0.3\onelineskip}%
    \fi%
    \ifthenelse{\equal{#1}{top}}{\vspace*{\fill}}{}%
  }
\else
  \newcommand{\iconographieimg}[2]{}
  \newcommand{\iconographietex}[1]{}
  \newcommand{\iconographiewrapfig}[4]{}
  \NewDocumentCommand{\iconographieinlineblock}{O{} m m}{}
  \NewDocumentCommand{\iconographiedouble}{O{bottom} O{5cm} m O{} m O{} O{}}{}
\fi

% ─── NOTES FINALES ────────────────────────────────────────────────────────────
% Compiler avec : lualatex "\def\AVECNOTES{}\input{...}"
\usepackage{endnotes}
\ifdefined\AVECNOTES
  \newcommand{\nf}[1]{\endnote{#1}}
\else
  \newcommand{\nf}[1]{}
\fi
% Supprimer le titre automatique du package (on gère le nôtre)
\renewcommand{\notesname}{}
% Style des notes : même fonte, corps légèrement réduit
\renewcommand{\enotesize}{\small}
% Format des notes en fin de livre : sans retrait, numéro en corps normal
\makeatletter
\renewcommand\enoteformat{%
  \rightskip\z@ \leftskip\z@ \parindent=0pt
  \leavevmode\theenmark.\enspace}
\makeatother
% Source dans les notes : saut de ligne avant le lien Wikipedia
\newcommand{\source}[1]{\newline\textit{Source~:}~{\upshape\def\_{\textunderscore\allowbreak}\def\%{\char"25\relax\allowbreak}#1}}
% Image seule dans une note (centrée sous le texte)
\newcommand{\nfimg}[2]{\newline\includegraphics[width=#1\linewidth]{#2}}
% Image côte à côte avec le texte : #1=texte, #2=frac image, #3=path, #4=url source
\newcommand{\nfimgblock}[4]{%
  \begin{minipage}[t]{0.45\linewidth}\small#1\end{minipage}%
  \hfill%
  \begin{minipage}[t]{#2\linewidth}\includegraphics[width=\linewidth]{#3}\end{minipage}%
  \source{#4}}
% Filet ornemental (séparateur de scène intra-chapitre)
\newcommand{\scenbreak}{%
  \vspace{\onelineskip}%
  {\centering\rule{0.4\textwidth}{0.4pt}\par}%
  \vspace{\onelineskip}%
}
% Titre de chapitre dans le bloc de notes
\newcommand{\chapnoteheader}[1]{{\normalfont\scshape #1}}
% Sous-chapitre intermédiaire : vide les notes courantes, ouvre un nouveau chapitre, repart à 1
% Usage : \startnewchapter        (titre auto = numéro romain du chapitre suivant)
%         \startnewchapter[Titre] (titre explicite)
% Actif uniquement si \CHAPITRESSEPARES est défini dans le document racine.
\ifdefined\CHAPITRESSEPARES
  \newcommand{\startnewchapter}[1][]{%
    \ifdefined\AVECNOTES
      \cleardoublepage
      \def\@snc@title{#1}%
      {\centering\normalfont\scshape\large
        Notes\ifx\@snc@title\empty\else~\textemdash~#1\fi\par}%
      \vspace{\onelineskip}%
      \begingroup\parindent=0pt
      \theendnotes
      \endgroup
    \fi
    \chapter{\ifx&#1&\Roman{chapter}\else#1\fi}%
    \beginchapternotes{\ifx&#1&\Roman{chapter}\else#1\fi}%
  }
\else
  \newcommand{\startnewchapter}[1][]{}%
\fi
% Notes par chapitre : remet le compteur à 0 pour que chaque acte reparte de 1
\newcommand{\beginchapternotes}[1]{%
  \ifdefined\AVECNOTES
    \setcounter{endnote}{0}%
    \ifdefined\NOTESPARTCHAPITRE\else
      \addtoendnotes{\protect\par\protect\medskip\protect\noindent\protect\chapnoteheader{#1}\protect\par}%
    \fi
  \fi
}

% ─── PAGE GEOMETRY (155×235 mm – Medium 1) ───────────────────────────────────
\setstocksize{235mm}{155mm}
\settrimmedsize{235mm}{155mm}{*}
\settrims{0mm}{0mm}
\setlrmarginsandblock{23mm}{18mm}{*} % inner/outer
\setulmarginsandblock{25mm}{28mm}{*} % top/bottom
\setheadfoot{11mm}{13mm}
\setheaderspaces{*}{7mm}{*}
\checkandfixthelayout

% ─── TYPOGRAPHY ───────────────────────────────────────────────────────────────
\raggedbottom
\renewcommand{\baselinestretch}{1.2}
\setlength{\parindent}{4mm}
\setlength{\parskip}{0pt}
\setlength{\afterchapskip}{0.4\onelineskip}

% ─── DIALOGUE ─────────────────────────────────────────────────────────────────
\newenvironment{dialogue}
  {\par\vspace{0.5\baselineskip}\setlength{\parindent}{0pt}}
  {\par\vspace{0.5\baselineskip}}


% ─── TABLE DES MATIÈRES ───────────────────────────────────────────────────────
\setcounter{tocdepth}{0}
\renewcommand{\chapternumberline}[1]{}           % suppress arabic number, keep Roman title only
\renewcommand{\cftchapterdotsep}{\cftdotsep}     % dot leaders
\renewcommand{\cftchapterfont}{\normalfont\scshape}     % chapter entry: small caps, not bold
\renewcommand{\cftchapterpagefont}{\normalfont\scshape} % page number: small caps, not bold
% TOC title — gallimard chapter style suppresses \printchaptertitle, so we override here
\renewcommand{\printtoctitle}[1]{{\centering\normalfont\scshape\large #1\par}}
\renewcommand{\aftertoctitle}{\par\vspace{\onelineskip}}
% \mytableofcontents : cleardoublepage + suppress self-reference entry
\newcommand{\mytableofcontents}{%
  \cleardoublepage
  \begingroup
  \small\renewcommand{\baselinestretch}{1.0}\selectfont
  \addtocontents{toc}{\protect\setcounter{tocdepth}{-1}}%
  \tableofcontents
  \addtocontents{toc}{\protect\setcounter{tocdepth}{0}}%
  \endgroup
}

% ─── CHAPTER STYLE ────────────────────────────────────────────────────────────
\makeatletter
\makechapterstyle{gallimard}{%
  \renewcommand{\chapnamefont}{\normalfont\scshape\centering}
  \renewcommand{\chapnumfont}{\normalfont\scshape\centering}
  \renewcommand{\chaptitlefont}{\normalfont\scshape\centering}
  \renewcommand{\printchaptername}{}
  \renewcommand{\chapternamenum}{}
  \renewcommand{\printchapternum}{%
    \chapnumfont\Roman{chapter}%
    \markboth{\Roman{chapter}}{\Roman{chapter}}}
  \renewcommand{\afterchapternum}{\par\vskip 0.4\onelineskip}
  \renewcommand{\printchaptertitle}[1]{}
}
\makeatother
\chapterstyle{gallimard}

% ─── RUNNING HEADERS : folio seul, sans auteur ni titre ───────────────────────
\makepagestyle{gallimard}
\makeevenhead{gallimard}{}{\footnotesize\scshape\leftmark}{}
\makeoddhead{gallimard}{}{\footnotesize\scshape\rightmark}{}
\makeevenfoot{gallimard}{}{\thepage}{}
\makeoddfoot{gallimard}{}{\thepage}{}
\makepsmarks{gallimard}{%
  \nouppercaseheads
}
\pagestyle{gallimard}
\aliaspagestyle{chapter}{empty}

% ─── EDITORIAL NOTE (Avertissement, Postfaces) ────────────────────────────────
% 10.95pt (~small), interligne ×1.1, sans alinéa, titre en petites capitales.
\newenvironment{editorialnote}[1]{%
  \begingroup
  \renewcommand{\baselinestretch}{1.1}%
  \small
  \setlength{\parindent}{0pt}%
  \setlength{\parskip}{\onelineskip}%
  {\centering\normalfont\scshape\large #1\par}%
  \vspace{\onelineskip}%
  \noindent\ignorespaces
}{%
  \par\endgroup
}
